Drop foot is the inability to lift the forefoot, most often from weakness of tibialis anterior (TA) after peroneal nerve injury, stroke or multiple sclerosis \cite{stewart2008}. It shows in two phases of gait. In swing the toes are not lifted, so toe clearance falls and the foot can catch the ground; patients compensate with extra hip and knee flexion (steppage) or circumduction. At heel strike the missing eccentric TA moment lets the foot slap onto the floor. The clinical answer is a passive ankle-foot orthosis (AFO), a spring that holds the foot near neutral but also resists the plantarflexion needed for push-off. Active AFOs change their impedance over the cycle: braking in early stance, transparent in late stance and lifting the foot in swing \cite{blaya2004}. To do that the device must know the gait phase, and most devices estimate it from a shank-mounted inertial measurement unit (IMU).

This project asks whether such a phase-based active AFO, triggered only by a realistic simulated shank gyroscope and limited by a sized actuator, does better than the best passive AFO in a neuromuscular walking simulation of drop foot. The model is the planar Geyer and Herr reflex walker \cite{geyer2010} in SCONE \cite{geijtenbeek2019}; drop foot is a weak right TA, and the model's compensation is the re-optimization of an asymmetric reflex controller with CMA-ES \cite{hansen2016}. A similar model was used by Dermksian et al.\ \cite{dermksian} at 33\% TA strength without a device. Here the device chain is built end to end: torque requirement, motor and series elastic actuator (SEA) sizing, torque control reduced to a fitted response model, a causal gait-event detector on a sampled, noisy, biased and delayed gyroscope, and two device controllers. The hypothesis was committed to the repository before any device comparison was run, and every condition is reported over three optimization seeds, with falls counted as results.
